\documentclass[11pt]{article}
\usepackage[margin=1.1in]{geometry}
\usepackage{times}
\usepackage[T1]{fontenc}
\usepackage{setspace}
\onehalfspacing
\usepackage{parskip}
\pagenumbering{gobble}

\title{We Take This Seriously}
\author{Flyxion}
\date{September 2026}

\begin{document}
\maketitle

Following an incident, the company takes the matter seriously.

Seriousness has a fixed set of components: concern, a commitment to learning, gratitude toward whoever surfaced the problem, and confidence that the existing safeguards were basically sound all along. What it does not require, and what never appears alongside it, is a specific account of what happened, a name for who is responsible, a remedy with a date attached, or a stated consequence if the date is missed. A recall notice has all four. A consent decree has all four. This sentence is built, with some craft, to have none of them, because a sentence with none of them cannot later be shown to have been broken.

It borrows the grammar of an apology without the machinery that makes an apology cost anything. "I apologize" only functions as a performative, an act accomplished by the saying of it, when the speaker remains answerable for it afterward against some standard outside their own say-so. Remove that and the sentence collapses back into what it always was underneath the grammar, a report on the speaker's present facial expression, filed with no expiration date and no one positioned to check it against anything at all. It persists across unrelated companies and unrelated harms not because they share a culture but because they share a target: not the people affected, who get nothing checkable from it, but the news cycle covering the story, which needs a quote and will accept one that promises nothing as readily as one that promises something specific and later gets held to it. The matter, having now been taken seriously, is no longer available for anyone to ask what, exactly, was done about it.

\end{document}
